\chapter{Study reference}
\small
\lstset{basicstyle=\ttfamily\small}
This is a supplementary reference, not an official exam appendix. Try mixed
review before opening it. These are the signatures/semantics used in our files.
\par\noindent\begin{minipage}{\linewidth}
\begin{lstlisting}
id      :: a -> a
const   :: a -> b -> a
(.)     :: (b -> c) -> (a -> b) -> a -> c
($)     :: (a -> b) -> a -> b
fst     :: (a,b) -> a
snd     :: (a,b) -> b
(,)     :: a -> b -> (a,b)
(:)     :: a -> [a] -> [a]
Just    :: a -> Maybe a
Nothing :: Maybe a
map     :: (a -> b) -> [a] -> [b]
filter  :: (a -> Bool) -> [a] -> [a]
concat  :: [[a]] -> [a]
length  :: [a] -> Int
null    :: [a] -> Bool
head    :: [a] -> a       -- partial on []
tail    :: [a] -> [a]     -- partial on []
even    :: Integral a => a -> Bool
show    :: Show a => a -> String
(+)     :: Num a => a -> a -> a
(==)    :: Eq a => a -> a -> Bool
max     :: Ord a => a -> a -> a
fmap, (<$>) :: Functor f => (a -> b) -> f a -> f b
pure    :: Applicative f => a -> f a
(<*>)   :: Applicative f => f (a -> b) -> f a -> f b
return  :: Monad m => a -> m a
(>>=)   :: Monad m => m a -> (a -> m b) -> m b
(>>)    :: Monad m => m a -> m b -> m b
mapM    :: Monad m => (a -> m b) -> [a] -> m [b]
sequence :: Monad m => [m a] -> m [a]
\end{lstlisting}
\end{minipage}\par
\par\noindent\begin{minipage}{\linewidth}
\begin{lstlisting}
newtype State s a = S {runState :: s -> (a,s)}
get       :: State s s
put       :: s -> State s ()
modify    :: (s -> s) -> State s ()
evalState :: State s a -> s -> a
execState :: State s a -> s -> s
\end{lstlisting}
\end{minipage}\par
\clearpage
\section*{State and verification reminders}
\code{runState} returns \code{(result, finalState)}. Function results cannot
normally be printed; apply them to an argument before comparing output.
\code{return} constructs an action; it does not exit a \code{do} block.
\code{Maybe} bind propagates \code{Nothing}; list bind concatenates each
branch's results in order. State bind passes the updated state to the next action.

\par\noindent\begin{minipage}{\linewidth}
For an invariant $I$, guard $B$, body $S$, initialization $A$, suffix $T$,
method precondition $P$, and postcondition $Q$, our partial-correctness obligations are
\[
 P\Rightarrow wp(A,I),\qquad
 I\land B\Rightarrow wp(S,I),\qquad
 I\land\neg B\Rightarrow wp(T,Q).
\]
Assignment substitutes in the postcondition; sequence works backward.
For conditionals, use guarded implications for both branches. These obligations
do not by themselves prove termination or all runtime-safety properties.

\end{minipage}\par
